\scanpage{37}\mspage{15 bis}

\[
\deleted{=(\alpha w)^{-1}(\alpha g\unsure{u|v})(\alpha gv)^{-1}(\alpha gfu)
 ((\alpha w)^{-1}(\alpha gfu))^{-1}=1}
\]

\textcircled{0} $M$ pseudocofiltrant i.e. $\forall u,v\quad\exists u',v'$ avec $uu'=vv'$

\textcircled{1} $\forall u$, $\exists u'$ tel que $uu'\sim\psi$ i.e. $\exists p,q\in\mathbb N$ tels que
\[
\varphi^p uu'=\varphi^q\psi
\]

\textcircled{2} $\exists a\xrightarrow{\lambda}e=a\psi$ i.e. $\exists\lambda$ tel que $\psi\lambda\sim1$ i.e.
\[
\exists\alpha,\beta\in\mathbb N\text{ tels que }\quad\boxed{\varphi^\alpha\psi\lambda=\varphi^\beta}
\]

[ \uline{NB} Si $\psi$ est \uline{minimal}, alors \textcircled{0} et \textcircled{1} sont
satisfaits, et \textcircled{2} est satisfait sss \deleted{\unsure{Question 3}}
$\exists\beta\in\mathbb N$ tel que $\varphi^\beta$ soit minimal \unsure{\}.|?|(.}. Et
pour un tel $\beta$, $\exists\lambda\in M$ (et $\unsure{=} $ $\lambda$ \unsure{minimal|minim.?})
tel que $\psi\lambda=\varphi^\beta$ ]

\textcircled{3} $\forall r,s\in\mathbb N$ avec $\varphi^r\psi=\varphi^s\psi$, on a $r=s$.

[ NB Si $\psi$ est minimal, alors la relation
\[
\mathscr R_0:\qquad u\overset{R_0}{\sim}v\quad\Longleftrightarrow\quad u\psi=v\psi
\]
est aussi la relation d'équivalence définie par
l'application
\[
\begin{aligned}
M&\longrightarrow \unsure{\mathrm{Ep}|\mathrm{E}p}(M_0)\\
u&\longmapsto u_{M_0}=(\lambda\longmapsto u\lambda:M_0\longrightarrow M_0),
\end{aligned}
\]
(puisque $u\psi=v\psi\Longrightarrow u\psi\deleted{\unsure{\beta}}\alpha=v\psi\deleted{\unsure{\beta}}\alpha\quad\forall\alpha\in M$,
or $\psi M=M_0$ i.e. l'ens des $\psi\alpha$ est l'ens des
$\lambda\in M_0$). Soit $\bar M\simeq M/R_0\subset\mathrm{Ep}(M_0)$ le
monoïde image. \deleted{Ainsi, le} C'est un monoïde
dans lequel on a : $\bar u\bar f=\bar v\bar f\Longrightarrow\bar u=\bar v$, et
il est pseudocofiltrant, comme image de
$M$. (Il n'est filtrant que s'il est réduit à $1$.)
Ainsi la relation $\varphi^r\psi=\varphi^s\psi$ dans \textcircled{3} équivaut
à $\bar\varphi^r=\bar\varphi^s$, où $\bar\varphi$ est l'image de $\varphi$ dans $\bar M$.
Si on a $r=s+h$, $h\geq1$, \deleted{on en conclut} \inserted{cela équivaut} à

% archive sheet number: 36
